\chapter{序論}

本章では，本研究が取り組む問題の背景と課題を述べ，本研究の目的と本論文の構成について説明する．

\section{研究背景}

2022 年の ChatGPT の公開以降，大規模言語モデル（Large Language Model，以下 LLM）の利用が急速に広がり，専門的な訓練を受けていない利用者が日常的に使う道具になっている．日本語環境でも Swallow \cite{fujii2024swallow} のような日本語を主対象としたモデルが公開され，評価基盤も整備されつつある \cite{llmjpeval2024}．

利用の広がりに伴い，LLM の出力をどこまで信じてよいかという問題が生じている．LLM は事実と異なる内容を生成することがあり，この現象は幻覚（hallucination）と呼ばれる \cite{ji2023hallucination}．ここで問題になるのは出力が誤っていること自体ではなく，誤った出力が正しい出力と区別のつかない語調で提示されることである．利用者は，どの回答に検証の手間を割くべきかを判断しにくい．

この問題への対処として，LLM に自身の回答への確信度（confidence）を言語で表明させる手法が研究されており，verbalized confidence と呼ばれる \cite{kadavath2022know}\cite{lin2022teaching}\cite{tian2023justask}．表明された確信度が有用であるためには，それが実際の正答率と対応している必要がある．確信度 0.9 と答えた問題群の正答率が実際に 9 割程度であれば，利用者は「確信度 0.9 以上の回答は採用し，それ未満は検証する」という判断ができる．この対応の度合いを較正（calibration）と呼び，その程度を測る代表的な指標が ECE（Expected Calibration Error）である \cite{guo2017calibration}．

\section{課題}

較正に関する研究の多くは英語のタスクを対象としている．多言語での較正を扱った研究として MlingConf \cite{xue2025mlingconf} があるものの，日本語のタスクに焦点を当てた検証は乏しく，日本語の利用者が確信度表示をどこまで信頼できるかという実用上の問いに，現状では答えられない．

また，確信度の引き出し方によって較正の結果が変わることが複数の研究で指摘されている \cite{tian2023justask}\cite{xiong2024llmuncertainty}．回答と確信度を同時に尋ねるか，回答を得てから改めて尋ねるか，数値で答えさせるか言語表現で答えさせるかといった違いが結果に影響する．方式の比較そのものは英語では行われているが，日本語のタスクで，同一のモデル・同一のデータ・同一の評価指標のもとに複数の方式を並べた報告は限られている．

さらに，評価指標の選択自体も論点である．ECE の値は確信度の分布に左右され，確信度が狭い範囲に偏っていると，全体の平均確信度と平均正答率の差に近い値をとりうる．本研究の予備実験でもこの状況が生じた（5.4 節）．

\section{本研究の目的と本論文の構成}

本研究は，日本語タスクにおいて，確信度の引き出し方の違いが LLM の較正に与える影響を明らかにし，較正を改善するプロンプト設計の指針を得ることを目的とする．具体的な問いは 3 つである．第 1 に，確信度の引き出し方式によって較正はどのように変化するか．第 2 に，観察された差は較正の水準の違いによるものか，正答率の違いに伴うものか．第 3 に，較正を改善するためにプロンプトをどのように設計すべきか．

第 1 の問いに答えるため，3 つの引き出し方式を設計し，同一条件で比較する評価基盤を構築した．第 2 の問いは，本論文の範囲では方式によって回答そのものが変わるため完全には分離できていない．第 3 の問いは未着手であり，今後取り組む．

以降，第 2 章では関連研究を，第 3 章では設計方針を，第 4 章では実装を述べ，第 5 章で予備実験の結果と考察を，第 6 章でまとめと今後の方針を示す．
